
\subsection{Controle das juntas do manipulador robótico}
\label{sec_controle_PD_juntas}

O \gls{mr} acoplado à espaçonave $\{B\}$ possui 3 graus de liberdade rotacionais, descritos pelas coordenadas generalizadas, por suas respectivas velocidades articulares $\dot{\vec\phi}$:

\begin{equation}
\vec\phi =
\begin{bmatrix}
\phi_1 \\ \phi_2 \\ \phi_3
\end{bmatrix}.
\end{equation}

As referências articulares são denotadas por:

\begin{equation}
\vec\phi_{\text{ref}} =
\begin{bmatrix}
\phi_{1,\text{ref}} \\
\phi_{2,\text{ref}} \\
\phi_{3,\text{ref}}
\end{bmatrix}.
\end{equation}

As quais são obtidas a partir da solução do problema de cinemática inversa, conforme a Seção~\eqref{sec_cinematica_inversa}, do efetuador para uma trajetória desejada de aproximação e alinhamento com a porta de acoplamento.
Define-se o erro articular e sua derivada como:

\begin{equation}
\vec e_\phi = \vec\phi - \vec\phi_{\text{ref}},
\qquad
\dot{\vec e}_\phi = \dot{\vec\phi} - \dot{\vec\phi}_{\text{ref}}.
\end{equation}

Em que as referências de velocidade $\dot{\vec\phi}_{\text{ref}}$ são nulas ou variam lentamente, de modo a caracterizar trajetórias suavizadas para o manipulador.
A lei de controle \gls{pd} em espaço de juntas é então escrita como:

\begin{equation}
\vec\tau_\phi
=
- K_{p,\phi}\,\vec e_\phi
- K_{d,\phi}\,\dot{\vec e}_\phi.
\label{eq_controle_PD_juntas}
\end{equation}

O termo $\vec\tau_\phi$ é o vetor de torques nas juntas,
e $K_{p,\phi}$, $K_{d,\phi}$ são matrizes com os ganhos de posição e velocidade, respectivamente.
Neste contexto, o \gls{mr} é controlado no espaço de juntas \cite{craig2005}, enquanto o vínculo com o espaço operacional (posição e orientação do efetuador em relação à porta de acoplamento) é realizado pela cinemática direta e inversa discutidas anteriormente.
